% \documentclass[tikz,border=10pt]{standalone}
% \usepackage[utf8]{inputenc}
% \usepackage{amsmath,amssymb}
% \usetikzlibrary{arrows.meta}
% 
% \begin{document}
% 
% % Estilos y colores locales al tikzpicture (nada global).
% % Solo TikZ básico: no hace falta pgfplots.
\begin{tikzpicture}[
  x=0.8cm, y=0.8cm,
  % ---------- Estilos ----------
  eje/.style={
    -{Stealth[length=2.2mm,width=1.6mm]},
    line width=0.7pt, color=ejecol
  },
  marca/.style={line width=0.6pt, color=ejecol},
  guia/.style={
    dashed, line width=0.6pt, color=guiacol
  },
  curva/.style={
    line width=1.5pt, color=curvacol, line cap=round
  },
  curvaflecha/.style={
    curva, -{Stealth[length=2.6mm,width=2mm]}
  },
  panel/.style={
    rounded corners=8pt, draw=curvacol!35, fill=curvacol!5,
    line width=0.8pt
  },
  nombre/.style={
    anchor=north, text=curvacol!85!black, font=\large\bfseries
  },
  % Ejes con marcas en +-1, dibujados a mano (no depende de librerías).
  % Uso: \pic {ejes};
  pics/ejes/.style={
    code={
      \draw[eje] (-3.4, 0) -- (3.6, 0) node[right, text=black] {$z$};
      \draw[eje] (0, -1.9) -- (0, 1.9);
      \draw[marca] (-0.1, 1) -- (0.1, 1);
      \draw[marca] (-0.1,-1) -- (0.1,-1);
      \node[left=3pt, font=\small, fill=white, inner sep=1pt, text=black] at (0, 1) {$1$};
      \node[left=3pt, font=\small, fill=white, inner sep=1pt, text=black] at (0,-1) {$-1$};
    }
  }
]
  % ---------- Colores (locales a este tikzpicture) ----------
  \definecolor{curvacol}{RGB}{25,100,190}
  \definecolor{ejecol}{RGB}{60,70,90}
  \definecolor{guiacol}{RGB}{140,150,170}

  % ================= Sigmoide =================
  \begin{scope}[shift={(0,0)}]
    \draw[panel] (-4.2,-3.7) rectangle (4.9,4.3);
    \draw[guia] (-3,1) -- (3,1);
    \pic {ejes};
    \draw[curva] plot[domain=-3:3, samples=80, smooth]
      (\x, {1/(1+exp(-\x))});
    \node[anchor=south] at (0.1,2.2) {\large $g(z)=\dfrac{1}{1+e^{-z}}$};
    \node[nombre] at (0,-2.3) {Sigmoide};
  \end{scope}

  % ================= Tangente hiperbólica =================
  \begin{scope}[shift={(10.6,0)}]
    \draw[panel] (-4.2,-3.7) rectangle (4.9,4.3);
    \draw[guia] (-3,1) -- (3,1);
    \draw[guia] (-3,-1) -- (3,-1);
    \pic {ejes};
    \draw[curva] plot[domain=-3:3, samples=80, smooth]
      (\x, {tanh(\x)});
    \node[anchor=south] at (0.1,2.2) {\large $g(z)=\dfrac{e^{z}-e^{-z}}{e^{z}+e^{-z}}$};
    \node[nombre] at (0,-2.3) {Tangente hiperb\'olica};
  \end{scope}

  % ================= ReLU =================
  \begin{scope}[shift={(0,-9)}]
    \draw[panel] (-4.2,-3.7) rectangle (4.9,4.3);
    \pic {ejes};
    \draw[curvaflecha] (-3,0) -- (0,0) -- (1.55,1.55);
    \node[anchor=south] at (0.1,2.2) {\large $g(z)=\max\{z,\,0\}$};
    \node[nombre] at (0,-2.3) {ReLU};
  \end{scope}

  % ================= alfa-ELU =================
  \begin{scope}[shift={(10.6,-9)}]
    \draw[panel] (-4.2,-3.7) rectangle (4.9,4.3);
    \draw[guia] (-3,-1) -- (3.2,-1)
      node[right, font=\small, text=black] {$-\alpha$};
    \pic {ejes};
    \draw[curvaflecha] plot[domain=-3:0, samples=60, smooth]
      (\x, {exp(\x)-1}) -- (1.55,1.55);
    \node[anchor=south] at (0.1,2.2)
      {\large $g(z)=\begin{cases}
         z & z\ge 0\\[2pt]
         \alpha\,(e^{z}-1) & z<0
       \end{cases}$};
    \node[nombre] at (0,-2.3) {$\alpha$-ELU};
  \end{scope}

\end{tikzpicture}

%\end{document}